\subsection{Zero-curvature equation and the Landau--Lifshitz equations of motion}
\label{sec:class5-eom}

We now determine the continuum Hamiltonian flow independently of the time Lax
matrix derived in Section~\ref{sec:class5-time-lax}.  This fixes the dynamical
identification left open in Section~\ref{sec:lattice-to-continuum-zc} and
provides a direct test of the corrected matrix $V_5$.

It is convenient to introduce the constant antisymmetric matrix
\begin{equation}
M_5
:=
\begin{pmatrix}
0&0&-1\\
0&0&-\ii\\
1&\ii&0
\end{pmatrix},
\qquad
M_5^{\mathsf T}=-M_5.
\label{eq:class5-M-definition}
\end{equation}
The lower symbol of a single Class 5 bond Hamiltonian, evaluated on the
coherent states at the neighbouring points $x$ and $x+\epsilon$, has the
long-wavelength expansion
\begin{equation}
 h_5^{\downarrow}
 \bigl(\boldsymbol S(x,t),\boldsymbol S(x+\epsilon,t)\bigr)
 =2+\epsilon^2\mathcal H_5(x,t)+\mathcal O(\epsilon^3),
\label{eq:class5-h-lower-expansion}
\end{equation}
where
\begin{equation}
\mathcal H_5
=
-\frac12\,\partial_x\boldsymbol S^{\mathsf T}\partial_x\boldsymbol S
+\frac{\alpha_5}{2}
\left(
\boldsymbol S^{\mathsf T}M_5\partial_x\boldsymbol S
+\partial_x S^+
\right).
\label{eq:class5-hamiltonian-density-with-boundary}
\end{equation}
The last term is a total spatial derivative.  With the periodic boundary
condition in \eqref{eq:continuum-spin-field}, the continuum Hamiltonian obtained
from the lattice Hamiltonian is therefore
\begin{align}
H^{\rm cl}_5
&:=
\lim_{\epsilon\to0}
\frac{
\bigl(H^{\rm lat}_5\bigr)^{\downarrow}-2N
}{\epsilon}
\nonumber\\
&=
\int_0^{\ell}dx\left[
-\frac12\,\partial_x\boldsymbol S^{\mathsf T}\partial_x\boldsymbol S
+\frac{\alpha_5}{2}
\boldsymbol S^{\mathsf T}M_5\partial_x\boldsymbol S
\right].
\label{eq:class5-continuum-hamiltonian}
\end{align}
The signs and the overall normalization in
\eqref{eq:class5-continuum-hamiltonian} follow from the lattice Hamiltonian
\eqref{eq:class5-lattice-hamiltonian}, the Heisenberg equation
\eqref{eq:Heisenberg-lattice-time}, and the continuum time convention
$t=\epsilon^2t_{\rm lat}$.

For the unit spin field normalized in \eqref{eq:continuum-spin-field}, we use
the continuum Poisson bracket
\begin{equation}
\{S^i(x),S^j(y)\}_5
=
2\,\varepsilon_{ijk}S^k(x)\,\delta(x-y),
\qquad
\varepsilon_{123}=1,
\label{eq:class5-poisson-bracket}
\end{equation}
where $\delta(x-y)$ is the Dirac delta function.  The corresponding Hamiltonian
flow is
\begin{equation}
\partial_t S^i(x)=\{S^i(x),H^{\rm cl}_5\}_5.
\label{eq:class5-hamiltonian-flow-definition}
\end{equation}
Using the bilinear cross product
$(\boldsymbol A\times\boldsymbol B)^i
=\varepsilon_{ijk}A^jB^k$, the equations of motion take the compact form
\begin{equation}
\partial_t\boldsymbol S
=
-2\,\boldsymbol S\times
\left(
\partial_x^2\boldsymbol S
+\alpha_5 M_5\partial_x\boldsymbol S
\right).
\label{eq:class5-vector-eom}
\end{equation}
The right-hand side is tangent to the unit-spin constraint, so
$\partial_t(\boldsymbol S^{\mathsf T}\boldsymbol S)=0$.

For comparison with the matrix entries of the Lax pair, the same equations
can be written in terms of the transverse components $S^{\pm}$ defined in
\eqref{eq:class5-spin-pm}.  Using
$\boldsymbol S^{\mathsf T}\partial_x\boldsymbol S=0$, one obtains
\begin{align}
\partial_t S^+
&=
2\ii\left(
S^+\partial_x^2S^3
-S^3\partial_x^2S^+
+\alpha_5 S^+\partial_xS^+
\right),
\label{eq:class5-eom-plus}\\
\partial_t S^-
&=
2\ii\left(
S^3\partial_x^2S^-
-S^-\partial_x^2S^3
+\alpha_5 S^+\partial_xS^-
\right),
\label{eq:class5-eom-minus}\\
\partial_t S^3
&=
\ii\left(
S^-\partial_x^2S^+
-S^+\partial_x^2S^-
+2\alpha_5 S^+\partial_xS^3
\right).
\label{eq:class5-eom-three}
\end{align}
A direct expansion of the lower symbol of the lattice Heisenberg derivative
then gives
\begin{equation}
\left(
\partial_{t_{\rm lat}}\widehat L_{5;a,n}
\right)^{\downarrow}
=
\epsilon^3\,\partial_tU_5(x_n,t;\lambda)
+\mathcal O(\epsilon^4),
\label{eq:class5-heisenberg-to-Ut}
\end{equation}
with $\partial_t\boldsymbol S$ given by
\eqref{eq:class5-vector-eom}.  Equation~\eqref{eq:class5-heisenberg-to-Ut}
therefore supplies the dynamical identification required on the left-hand
side of the continuum equation \eqref{eq:continuum-zc-with-C}.

To establish the off-shell relation between the curvature and the equations
of motion, we treat the time derivatives as independent variables and define
the three equation-of-motion residuals
\begin{align}
E_+
&:=
\partial_t S^+
-2\ii\left(
S^+\partial_x^2S^3
-S^3\partial_x^2S^+
+\alpha_5 S^+\partial_xS^+
\right),
\label{eq:class5-residual-plus}\\
E_-
&:=
\partial_t S^-
-2\ii\left(
S^3\partial_x^2S^-
-S^-\partial_x^2S^3
+\alpha_5 S^+\partial_xS^-
\right),
\label{eq:class5-residual-minus}\\
E_3
&:=
\partial_t S^3
-\ii\left(
S^-\partial_x^2S^+
-S^+\partial_x^2S^-
+2\alpha_5 S^+\partial_xS^3
\right).
\label{eq:class5-residual-three}
\end{align}
The curvature of the matrices $U_5$ and $V_5$ in
\eqref{eq:class5-spatial-Lax} and \eqref{eq:class5-time-Lax-compact} is
\begin{equation}
F^{(5)}_{tx}
:=
\partial_tU_5-\partial_xV_5+[U_5,V_5].
\label{eq:class5-curvature-definition}
\end{equation}
Using only the unit-spin constraint and its spatial derivatives, the curvature
factorizes as
\begin{equation}
F^{(5)}_{tx}
=
\frac{1}{4\lambda}
\begin{pmatrix}
E_3+2\alpha_5\lambda E_+
&
E_- -2\alpha_5\lambda E_3\\
E_+
&
-E_3
\end{pmatrix}.
\label{eq:class5-offshell-factorization}
\end{equation}
For $\lambda\neq0$, the residuals are recovered from the curvature entries by
\begin{align}
E_+&=4\lambda\bigl(F^{(5)}_{tx}\bigr)_{21},
\qquad
E_3=-4\lambda\bigl(F^{(5)}_{tx}\bigr)_{22},
\nonumber\\
E_-&=4\lambda\left(
\bigl(F^{(5)}_{tx}\bigr)_{12}
-2\alpha_5\lambda\bigl(F^{(5)}_{tx}\bigr)_{22}
\right).
\label{eq:class5-residuals-from-curvature}
\end{align}
Hence the zero-curvature condition $F^{(5)}_{tx}=0$ is equivalent to the
three continuum equations of motion
\eqref{eq:class5-eom-plus}--\eqref{eq:class5-eom-three}.

To compare with the direct continuum construction of
Ref.~\cite{deLeeuwFontanellaNietoGarcia2026}, we translate our conventions by
\begin{equation}
S^2\mapsto -S^2,
\qquad
\alpha_5\mapsto -\alpha_5,
\qquad
t\mapsto 2t.
\label{eq:class5-reference-convention-change}
\end{equation}
With these substitutions, Eq.~\eqref{eq:class5-vector-eom} takes the form of
the Class 5 continuum equation used in
Ref.~\cite{deLeeuwFontanellaNietoGarcia2026}.  The reflection of the second
spin axis conjugates $M_5$ by $\operatorname{diag}(1,-1,1)$, yielding the
deformation matrix used there.
